\section{Finite rational tests for integer solvability}
\label{sec:reduction}

We construct the finite rational tests announced in the introduction and
prove that success of every test forces an ordinary integer zero.  The
construction has an ordinary integer model.  Conversely, compactness will
embed a model of all the tests in an elementary extension of \(\Q\),
where arithmetic restrictions on its ring values will force an integer
solution.  The later sections prove the arithmetic inputs used here.

The elliptic point data first give an additive assignment of integer
indices.  A local comparison will recover those indices from ring values
and determine the two local test conditions.

\subsection{Elliptic indices and the local comparison}
\label{red:local-comparison-section}

Fix the number field, involution, and elliptic curve
\begin{equation}\label{ell:data}
 F=\Q(\sqrt2),\qquad d=75-53\sqrt2,\qquad
 E:\ y^2=x^3-d^2x,
\end{equation}
where \(\sigma\) is the nontrivial automorphism of \(F/\Q\).
The following result, proved in Section~\ref{ell:section}, supplies the
integer indices.

\begin{proposition}\label{ell:rank}
The group $E(F)$ has rank one.  Consequently there exist a positive integer
$m$ and a nontorsion point $P\in E(F)$ such that
\[
                         mE(F)=\Z P.
\]
We fix such $m$ and $P$ once and for all.
\end{proposition}

We now describe the point data that the recursive theory will impose.
Let \(R\) be an integral domain
of characteristic zero for which \(\sqrt2\notin\Frac(R)\), and write
\[
 R_F=R[\sqrt2],\qquad F_R=\Frac(R)(\sqrt2).
\]
For each \(a\in R\), consider a point \(T_a\in E(F_R)\) subject to
\begin{equation}\label{red:elliptic-axioms}
\begin{gathered}
 T_0=O,\qquad T_1=P,\qquad
 T_{a+a'}=T_a+T_{a'}\quad(a,a'\in R),\\
 T_a\in mE(F_R)\quad(a\in R),\\
 a\ne0\Longrightarrow T_a\text{ is finite and }x(T_a)y(T_a)\ne0.
\end{gathered}
\end{equation}
The theory will supply six unary ring-valued functions
\(x_0,x_1,x_2,y_0,y_1,y_2\). For \(a\ne0\), they are required to
satisfy \(x_2(a)y_2(a)\ne0\) and
\begin{equation}\label{red:coordinate-functions}
 x(T_a)=\frac{x_0(a)+x_1(a)\sqrt2}{x_2(a)},\qquad
 y(T_a)=\frac{y_0(a)+y_1(a)\sqrt2}{y_2(a)}.
\end{equation}
Their values at \(a=0\) are unrestricted, since \(T_0=O\).
In the integer model we will take \(T_a=aP\). Representing the
coordinates by ring elements will also let a uniform height bound on
\(R\) control their heights.

For the arithmetic analysis, consider an elementary extension of the
ordinary many-sorted structure containing \(\Q,\Z,\R,F\), the prime and
valuation relations, the height functions, the elliptic multiple maps,
and the finite-product functions. We call elements of the original
structure \emph{ordinary}; starred sorts satisfy the same first-order
statements with ordinary parameters. For the moment, suppose that \(R\)
embeds in the rational sort \({}^*\Q\) of such an extension. Its
fraction-field interpretation embeds \(F_R\) in \({}^*F\). A tilde
denotes the image of a ring element, field element, or point. Later
compactness will supply an extension and an embedded ring with these
properties.

The identity \(mE(F)=\Z P\) transfers to this extension: every point
of \(mE({}^*F)\) has a unique expression \(nP\), with
\(n\in{}^*\Z\). Define
\begin{equation}\label{red:eta}
                         \widetilde T_a=\eta(a)P.
\end{equation}
The group-law axioms give
\begin{equation}\label{red:additivity}
       \eta(a+a')=\eta(a)+\eta(a'),\qquad \eta(1)=1.
\end{equation}
If \(\eta(a)=0\), then \(T_a=O\), and the finite-point axiom forces
\(a=0\). Thus \(\eta:R\to{}^*\Z\) is injective and additive.
Conjugating \eqref{red:eta} gives
\(\widetilde T_a^\sigma=\eta(a)P^\sigma\).
We make no assumption about \(\eta(aa')\).

For nonzero \(h\in R\), the index ratio of \(T_{ha}\) and \(T_h\) is
\(\eta(ha)/\eta(h)\), which need not equal \(\eta(a)\).  We will
recover this ratio locally from point parameters, and use a second
comparison to identify it with \(\eta(a)\) to the required precision.

For a finite point with \(y\ne0\), put \(z=-x/y\), and set \(z(O)=0\).
Every nonzero multiple of \(P\) avoids the two-torsion points, so
\(z(nP)\) is defined for every integer \(n\).  Write
\[
 \ell(Z)=Z+\sum_{t\ge2}c_tZ^t\in F[[Z]],\qquad
 \ell_K(Z)=Z+\sum_{2\le t\le K}c_tZ^t
\]
for the formal logarithm in this parameter and its truncation through
degree \(K\).  The truncations are effectively computable from the
fixed equation.  Section~\ref{ell:section} proves the following uniform
statements about integer multiples of the fixed \(F\)-rational points
\(P,P^\sigma\).

\begin{proposition}\label{ell:formal}
There is a fixed finite set of rational primes $S_E$ with the following
properties.  Let $p\notin S_E$, let $v$ be any place of $F$ above $p$,
normalized by $v(p)=1$, and let $n,n'\in\Z$, with $n'\ne0$.  Assume
$nP,n'P$ reduce to $O$ at $v$.  Then
\begin{equation}\label{ell:ratio}
           v\!\left(\frac{z(nP)}{z(n'P)}\right)=v_p(n/n').
\end{equation}
For every integer $K\ge1$, if in addition $p>K+1$ and $v_p(n/n')\ge0$,
then $\ell_K(z(n'P))\ne0$ and
\begin{equation}\label{ell:truncation}
 v\!\left(\frac{\ell_K(z(nP))}{\ell_K(z(n'P))}-\frac n{n'}\right)
                         \ge K.
\end{equation}
Both assertions also hold at the same place $v$ after replacing
$E,P,\ell_K$ by $E^\sigma,P^\sigma,\ell_K^\sigma$ and imposing the
corresponding reduction-to-$O$ hypotheses on the conjugate points.
Here and below the valuation of zero is $+\infty$, so $n=0$ is included.
\end{proposition}

The two estimates provide different kinds of local information. Equation
\eqref{ell:ratio} determines the valuation of an index ratio from the
point parameters without assuming that the ratio is integral. Once an
index ratio is known to be integral, \eqref{ell:truncation} approximates
it by a ratio of finite polynomial evaluations. The next lemma combines both forms of control.

Reduction to \(O\) has its usual valuation meaning, transferred to the
elementary extension.
Only the stated assertions about integer multiples of the fixed
\(F\)-rational points \(P,P^\sigma\) and evaluations of the finite
polynomials \(\ell_K,\ell_K^\sigma\) are transferred; no infinite
formal series is evaluated in the extension.

\begin{samepage}
\begin{lemma}[Local comparison]\label{red:local-comparison}
Let \(K\ge1\) be an ordinary integer. Let \(q\) be a prime of the
ambient integer sort \({}^*\Z\), possibly nonordinary, with
\(q\notin S_E\) and \(q>K+1\), and let \(v_*\) be a place of
\({}^*F\) above \(q\), normalized by \(v_*(q)=1\).
Suppose there are nonzero \(h,j_0\in R\), chosen independently of
\(a\), such that for every \(a\in R\):
\begin{enumerate}[label=(\roman*),leftmargin=*]
\item The points \(\widetilde T_h,\widetilde T_{ha}\) reduce to \(O\)
at \(v_*\), the value \(\ell_K(z(\widetilde T_h))\) is nonzero, and
\[
 v_*\!\left(
 \frac{\ell_K(z(\widetilde T_{ha}))}{\ell_K(z(\widetilde T_h))}
                   -\widetilde a\right)\ge K.
\]
\item The points \(\widetilde T_{j_0}^{\sigma}\) and
\(\widetilde T_{(h-4)a}^{\sigma}\) reduce to \(O\) on the conjugate
curve at the same place \(v_*\), and
\[
 v_*\!\left(
 \frac{z(\widetilde T_{(h-4)a}^{\sigma})}
      {z(\widetilde T_{j_0}^{\sigma})}\right)\ge K.
\]
\end{enumerate}
Then every \(\widetilde a\) is integral at \(q\), and
\[
                       v_q(\widetilde a-\eta(a))\ge K
                       \qquad(a\in R).
\]
Here \(z(O)=0\). The denominator in (ii) is nonzero because
\(j_0\ne0\) and the nonzero points in \eqref{red:elliptic-axioms}
have nonzero \(x\)- and \(y\)-coordinates.
\end{lemma}
\end{samepage}

\begin{proof}
Transfer the formal-parameter valuation equality of
Proposition~\ref{ell:formal} to the conjugate points at \(v_*\).
The comparison in (ii) gives
\[
 v_q\!\left(\frac{\eta((h-4)a)}{\eta(j_0)}\right)\ge K.
\]
The denominator is a nonzero ambient integer, so its \(q\)-valuation
is nonnegative. Consequently
\begin{equation}\label{red:alternate-index}
                         v_q(\eta((h-4)a))\ge K.
\end{equation}
Put \(V_a=\eta(ha)\). Additivity applied to
\((h-4)a=ha-4a\), including the case \(a=1\), gives
\[
              V_a\equiv4\eta(a)\pmod{q^K},\qquad
              V_1\equiv4\pmod{q^K}.
\]
Since \(q\ne2\), \(V_1\) is a \(q\)-unit. The ambient integers
\(V_a,\eta(a)\) are \(q\)-integral, and hence
\begin{equation}\label{red:actual-slope}
       V_a/V_1\text{ is \(q\)-integral},\qquad
       V_a/V_1\equiv\eta(a)\pmod{q^K}.
\end{equation}
The primary points have indices \(V_1\) and \(V_a\). They reduce to
\(O\), the denominator index is nonzero, and their ratio is now known
to be integral. The truncated-logarithm assertion of
Proposition~\ref{ell:formal} therefore gives
\[
 v_*\!\left(
 \frac{\ell_K(z(\widetilde T_{ha}))}{\ell_K(z(\widetilde T_h))}
                         -\frac{V_a}{V_1}\right)\ge K.
\]
The comparison in (i) relates this same logarithm ratio to
\(\widetilde a\). Combining it with the last display and
\eqref{red:actual-slope} proves the required congruence, since
\(v_*\) restricts to \(v_q\) on \({}^*\Q\). It also proves
\(q\)-integrality of \(\widetilde a\). Zero numerator indices are
covered by the convention \(v_q(0)=+\infty\).

\end{proof}

The same \(h\) must serve every \(a\), because the calculation uses
both \(a\) and \(1\).  We will impose ring-language versions of these
comparisons and obtain their valuation hypotheses at primes detected by a
height estimate.

\subsection{Contact primes and ordinary integer witnesses}
\label{red:arithmetic-inputs}

Throughout the paper, \(\hgt\) denotes the absolute logarithmic Weil
height.  The next input identifies primes whose product controls the height
of a rational point.  It is proved in Section~\ref{hei:section}.

\begin{theorem}[Five-point height estimate]\label{hei:height}
There exist five distinct points $\mathcal C\subset\mathbf P^1(\Q)$,
including $\infty$, a finite set $S$ of rational primes, and constants
$H,c>0$ with the following property.  For each $b\in\mathcal C$, choose
a primitive integral linear form $L_b(U,V)$ vanishing at $b$, with
$L_\infty(U,V)=V$. For
$s=(u:v)\in\mathbf P^1(\Q)\setminus\mathcal C$, let $M_S(s)$ be the
product, taken without repetition, of primes $q\notin S$ for which
at least one of the integers
\begin{equation}\label{hei:radical}
 v_q(L_b(u,v))-\min\{v_q(u),v_q(v)\},\qquad b\in\mathcal C,
\end{equation}
is odd. Here $v_q(q)=1$ and $v_q(0)=+\infty$. Then
\begin{equation}\label{hei:height-bound}
                    \hgt(s)\le H M_S(s)^c.
\end{equation}
\end{theorem}

Multiplying $(u,v)$ by a nonzero rational scalar adds the same valuation
to both terms in~\eqref{hei:radical}, so each contact order is independent
of the chosen homogeneous coordinates.  It is nonnegative after scaling
the coordinates to be integral and primitive at $q$.  With fixed primitive
integral coordinates, the nonzero integers $L_b(u,v)$ have only finitely
many prime divisors, so the product defining $M_S(s)$ is finite.



For a positive integer \(K\) and a nonzero integer \(e\) satisfying
\(M_S(s)^K\mid e\), the estimate gives
\[
                         \hgt(s)\le H|e|^{c/K}.
\]
In the reduction, the same polynomial error will play the role of \(e\)
for representations of every ring element.  This is why the exponent
\(c\) must be fixed, although its precise value is not needed.
The following immediate consequence permits the further fixed exclusions
needed by the local tests.

\begin{corollary}\label{hei:enlarge-exceptions}
If $S'\supset S$ is any fixed finite set of primes, then
\[
 \hgt(s)\le H\left(\prod_{q\in S'\setminus S}q\right)^c
                      M_{S'}(s)^c
 \qquad(s\in\mathbf P^1(\Q)\setminus\mathcal C).
\]
In particular enlarging the exceptional set changes the constant
but leaves the exponent unchanged.
\end{corollary}

\begin{proof}
Every prime counted by $M_S(s)$ is either counted by $M_{S'}(s)$ or
belongs to $S'\setminus S$.  Thus
$M_S(s)\le(\prod_{q\in S'\setminus S}q)M_{S'}(s)$, and
Theorem~\ref{hei:height} applies.
\end{proof}

To apply this estimate uniformly, the theory will write each ring element
as a sum of three fractions and impose local conditions on factors of their
contact forms.  The following ordinary integer result supplies witnesses
for those axioms.  It is a direct consequence of the prime-pattern lemma
and the Chinese remainder calculation proved in Section~\ref{pri:section}.

\begin{corollary}[Ordinary contact representations]\label{pri:integer-coverage}
Let \(\mathcal K\subset\Z_{>0}\) be the finite multiplier set supplied
by Lemma~\ref{pri:pattern} in Section~\ref{pri:section}.  It has the
following property.  For every integer \(K\ge1\), every \(A\in\Z\), and every
real lower bound \(Y_0>0\), there exist integers \(u_1,u_2,v\), with \(v\ne0\),
such that, on setting \(u_3=Av-u_1-u_2\), the following hold.
\begin{enumerate}[label=(\arabic*),leftmargin=*]
\item \textup{Representation and primes.} Each of the thirteen numbers
\begin{equation}\label{pri:coverage-list}
 v,\qquad L_b(u_i,v)
       \quad(1\le i\le3,\ b\in\mathcal C\setminus\{\infty\})
\end{equation}
equals \(\varepsilon k r\), where \(\varepsilon\in\{1,-1\}\),
\(k\in\mathcal K\), and \(r>Y_0\) is a positive rational prime.  Every
such \(r\) splits in \(F\), is a prime of good reduction for both
\(E\) and \(E^\sigma\), lies outside \(S_E\), and satisfies \(r>K+1\).

\item \textup{Witnesses at each place.} For each occurrence of a prime
\(r\) in this list and each of the two places \(v_r\) of \(F\) above it,
there are nonzero integers \(h,j_0\), chosen independently of \(a\in\Z\),
such that every \(a\in\Z\) satisfies conditions (i) and (ii) of
Lemma~\ref{red:local-comparison} in the ordinary setting: take \(R=\Z\),
its identity embedding in the ordinary structure, \(T_a=aP\), \(q=r\),
and \(v_*=v_r\), normalized by \(v_r(r)=1\).  In this specialization
the tildes disappear.  The choices of \(h,j_0\) may differ at different
occurrences and at the two places.
\end{enumerate}
The set \(\mathcal K\) is independent of \(K,A,Y_0\).
\end{corollary}

The two places in part (2) are the two residue branches above a split
prime.  At either one, conditions (i) and (ii) use the two conjugate curves
at that same place.  If \(\mathfrak p\) is the prime ideal of
\(\Z[\sqrt2]\) corresponding to \(v_r\), its valuation ring is
\((\Z[\sqrt2])_{\mathfrak p}\), with uniformizer \(r\).  Thus the two
valuation inequalities in part (2) are exactly membership in
\(r^K(\Z[\sqrt2])_{\mathfrak p}\); for a finite indicated point,
reduction to \(O\) is equivalent to
\(1/x\in r(\Z[\sqrt2])_{\mathfrak p}\).  We now encode these ordinary
witnesses by formulas that also make sense for a general ring.

\subsection{A recursive theory with an ordinary integer model}
\label{red:theory}

Let \(\mathcal T\) be a theory in a countable recursive expansion of
the language of rings. Its ring \(R\) is required to be an integral
domain of characteristic zero with \(\sqrt2\notin\Frac(R)\).
Add the six unary ring-valued functions above. Interpret \(T_0\) as
\(O\) and, for \(a\ne0\), require the two fractions in
\eqref{red:coordinate-functions} to define the coordinates of
\(T_a\in E(F_R)\), with both denominators nonzero. Impose the point
conditions \eqref{red:elliptic-axioms}.

Use the five points \(\mathcal C\), primitive contact forms \(L_b\), and
finite multiplier set \(\mathcal K\) of the preceding inputs. For each
ordinary integer \(K\ge1\), add the following representation axiom.
For every \(A\in R\) there exist \(u_1,u_2,v\in R\), with \(v\ne0\),
such that, on setting
\[
                            u_3=Av-u_1-u_2,
\]
each member of
\begin{equation}\label{red:contact-list}
 v,\qquad L_b(u_i,v)\quad
       (1\le i\le3,\ b\in\mathcal C\setminus\{\infty\})
\end{equation}
equals \(\varepsilon k r\), where
\(\varepsilon\in\{1,-1\}\), \(k\in\mathcal K\), and \(r\in R\)
has the following properties. The choices of \(r\) and its auxiliary
data are made separately for the different members of the list.

First require \(r\ne0\), \(rR\) proper, and \(R/rR\) a field.
Require \(B_r\in R\) satisfying
\[
 B_r^2\equiv2\pmod{rR},\qquad
 2B_r\text{ invertible in }R/rR.
\]
These conditions define two maximal ideals of \(R_F\):
\begin{equation}\label{red:branch-ideals}
 \mathfrak p_\pm
 =\ker\bigl(R_F\longrightarrow R/rR,\ \sqrt2\longmapsto\pm B_r\bigr).
\end{equation}
They are the two residue branches to be tested. In the integer model
they correspond to the two places of \(F\) above a supplied split
prime \(r\). In a general model they are defined by the quotient field,
independently of an ambient valuation.

For each \(\mathfrak p\in\{\mathfrak p_+,\mathfrak p_-\}\), require
nonzero \(h,j_0\in R\), chosen at that branch and independently of
\(a\), such that the following hold for every \(a\in R\). A point is
called \emph{formal for this localization} if it is \(O\), or if
\(1/x\in r(R_F)_{\mathfrak p}\).
\begin{enumerate}[label=(\roman*),leftmargin=*]
\item The points \(T_h,T_{ha}\) are formal for this localization,
\(\ell_K(z(T_h))\ne0\), and
\begin{equation}\label{red:primary-test}
 \frac{\ell_K(z(T_{ha}))}{\ell_K(z(T_h))}-a
                            \in r^K(R_F)_{\mathfrak p}.
\end{equation}
\item The points \(T_{j_0}^{\sigma},T_{(h-4)a}^{\sigma}\) are formal
for the conjugate curve at the same localization, and
\begin{equation}\label{red:alternate-test}
 \frac{z(T_{(h-4)a}^{\sigma})}{z(T_{j_0}^{\sigma})}
                            \in r^K(R_F)_{\mathfrak p}.
\end{equation}
\end{enumerate}
Here \(z(O)=0\), and the denominator in
\eqref{red:alternate-test} is nonzero by \(j_0\ne0\) and
\eqref{red:elliptic-axioms}. The pairs \(h,j_0\) may differ between
the two branches. The scheme contains every ordinary \(K\); the
termination proof will later select a single precision after \(f\)
is fixed.

All these conditions have recursive first-order expressions in \(R\).
The absence of \(\sqrt2\) from \(\Frac(R)\) is expressed by
\[
       \forall a,b\in R\quad b\ne0\Longrightarrow a^2\ne2b^2.
\]
Elements of \(R_F\) are pairs of ring elements, and elements of
\(F_R\) are fractions \((A+B\sqrt2)/D\), with \(A,B,D\in R\)
and \(D\ne0\). Equations with fixed coefficients in \(F\) can
therefore be separated into two coordinates and cleared of denominators.
The elliptic group law, including the identity and exceptional addition
cases, is given by finitely many field equations and inequations.
Membership in \(mE(F_R)\) is expressed by existential coordinates for
a point whose \(m\)-multiple is the prescribed point. These conditions
are ring formulas under the same interpretation.

Saying that \(R/rR\) is a field means that \(1\notin rR\) and every
element outside \(rR\) has an inverse modulo \(rR\). For \(D'\in F_R\),
the assertion
\begin{equation}\label{red:localization-definition}
                         D'\in r^i(R_F)_{\mathfrak p_\pm}
\end{equation}
means that \(D'=r^iX/Y\) for some \(X,Y\in R_F\) such that the
evaluation of \(Y\) under \(\sqrt2\mapsto\pm B_r\) is nonzero modulo
\(rR\), hence a unit in that field. These representations and inverse conditions
are expressed by ring quantifiers. The exponent \(i\) is fixed in
each axiom; neither variable exponentiation nor a valuation function
is part of \(\mathcal T\).

\begin{lemma}\label{red:standard-model}
The theory \(\mathcal T\) is recursive and has a model whose ring is
\(\Z\). In this model \(T_a=aP\) for every integer \(a\).
\end{lemma}

\begin{proof}
The ring and elliptic axioms hold for \(\Z\) with \(T_a=aP\).
For nonzero \(a\), the coordinates lie in \(F\) and admit integer
numerator and denominator representatives; choose arbitrary integer
values for the six coordinate functions at \(a=0\).
Nonzero multiples of the nontorsion point \(P\)
avoid \(O\) and the two-torsion points, so their \(x\)- and
\(y\)-coordinates are nonzero.

Fix \(K\ge1\) and \(A\in\Z\).  Corollary~\ref{pri:integer-coverage}
supplies \eqref{red:contact-list} with the required fixed multipliers and
positive split primes.  For each such \(r\), the quotient \(\Z/r\Z\)
is a field.  Splitness supplies a root \(B_r\) of \(2\) modulo \(r\),
and \(2B_r\) is invertible because \(r\) is odd.  The ideals
\eqref{red:branch-ideals} correspond to the two places above \(r\).
At each one, the corollary supplies nonzero \(h,j_0\), fixed for all
integers \(a\), with the required formal points and both valuation
comparisons.  Its localization interpretation gives exactly
\eqref{red:primary-test} and \eqref{red:alternate-test}.  Thus every
representation axiom holds.

The coefficients of \(\ell_K\) are computable from the fixed equation,
and the interpretations and axiom schemes above are effective. The
finite sets, fixed integers, rational coefficients, and coordinates
of \(P\) are finitely many exact constants that can be hardcoded in
the theory. Their existence suffices for the existence of the decision
procedure constructed below; no effective search for those fixed
choices is required. No list of all true arithmetic sentences is used.
\end{proof}

\subsection{Finite rational tests and compactness}
\label{red:finite-tests}

The finite tests will also impose a positive existential condition on
every ring value.  This input holds on all integers and controls the parity
of possible poles; Section~\ref{par:section} constructs it.

\begin{proposition}\label{par:pole-parity}
There are a positive existential formula $\Phi(b)$ in the language of
rings, with fixed rational coefficients, and a finite set $S_1$ of
rational primes such that:
\begin{enumerate}
\item $\Q\models\Phi(b)$ for every $b\in\Z$;
\item if $b\in\Q^\times$ and $\Q\models\Phi(b)$, then, for every
$q\notin S_1$, the inequality $v_q(b)<0$ implies $v_q(b)\in2\Z$.
\end{enumerate}
\end{proposition}

The second property does not rule out even pole orders. Its later use rests on a
simple product observation. If $q\notin S_1$ and nonzero $a,c\in\Q$
satisfy $\Phi(a)$ and $\Phi(ac)$, with both $v_q(a)$ and $v_q(ac)$
negative, then
\[
                  v_q(c)=v_q(ac)-v_q(a)\in2\Z.
\]
Section~\ref{red:valuation} applies this observation when $\Phi$ holds
on every element of the ring supplied by the finite tests.

An integer zero can therefore supply rational witnesses for all the
conditions imposed below.  In the converse direction, pole parity will
connect a contact prime in the ambient integer structure to one of the
tested localizations of the embedded ring.

Assume that an algorithm decides rational solvability of integral
polynomials. Given $f\in\Z[X_1,\ldots,X_n]$, the constant case can be
settled immediately, so assume $f$ is nonconstant. Add constants
$\mathbf a=(a_1,\ldots,a_n)$ and the equation $f(\mathbf a)=0$ to
$\mathcal T$, and Skolemize the resulting theory. This is effective for
a recursive first-order theory: after prenex conversion, each existential
quantifier is replaced by a function of the preceding universal variables.
All resulting function symbols are ring-valued. The Skolemized axioms
are universal sentences with quantifier-free matrices.

A ground term is a term with no variables, built from the named constants
using the ring operations and added functions in this language. A compatible
assignment of values to all ground terms will interpret the added functions
on the ring of term values. Enumerate these terms, associate a rational
variable to each, and enumerate the following constraints:
\begin{enumerate}[label=(\alph*),leftmargin=*]
\item the constants $0,1$ and the ring operations have their usual values;
\item equal input tuples have equal outputs for each function symbol;
\item every ground instance of every universal axiom holds;
\item the value of every ground term satisfies the positive existential
formula $\Phi$ from Proposition~\ref{par:pole-parity}.
\end{enumerate}
The witnesses for $\Phi$ in (d) are auxiliary rational variables.
They are not values of additional functions of the ring and do not create
new ground terms subject to (d). The same distinction applies to the
inverse variables introduced only when translating a finite test to
polynomial equations; they are not added to the term-generated ring.

Every finite initial part of this list is decidable by the assumed
rational-solvability algorithm. Indeed, after introducing the witnesses
for $\Phi$, the finite collection is an existential Boolean combination
of polynomial equalities and inequations. An inequation $g\ne0$ becomes
$gy=1$ with a new variable; a finite Boolean formula can be put into
disjunctive normal form; and a conjunction $g_1=\cdots=g_s=0$ is
equivalent over $\Q$ to $\sum_i g_i^2=0$. Clearing fixed rational
denominators gives finitely many integral polynomial queries.

If any test fails, report that $f$ has no integer zero. This answer is
sound: an integer zero extends the model of Lemma~\ref{red:standard-model}
and admits Skolem functions; all its ground-term values are integers
and hence satisfy $\Phi$.

\begin{lemma}\label{red:compactness}
If every finite rational test succeeds, there exist an ambient
elementary extension of the many-sorted structure specified in
\S\ref{red:local-comparison-section}, a
model \(R\) of \(\mathcal T\), a tuple \(\mathbf a\in R^n\) with
\(f(\mathbf a)=0\), and an embedding \(R\hookrightarrow{}^*\Q\)
into its rational sort such that every element of the image satisfies
\(\Phi\).
\end{lemma}

\begin{proof}
Take the ordinary many-sorted structure specified in
\S\ref{red:local-comparison-section}, with all the indicated functions
and relations. Expand this ambient structure by constants naming all
ordinary elements in their respective sorts. Its elementary diagram
consists of all first-order sentences in its language that are true in
the expanded structure. Use this diagram and add constants
assigning a value to every ground term of the Skolemized ring language,
the ring and functionality constraints, all ground universal instances,
and $\Phi$ for each term value. Every finite subset has a
model: its finitely many constraints are among a successful rational
test, and they can be interpreted in the ordinary structure. Finitely
many sentences of its elementary diagram are already true there.
The compactness theorem therefore supplies a simultaneous assignment
in an elementary extension; see \cite{Marker2002}.

Let $R$ be the set of values of the ground terms. Constraint (a) makes
it a subring of ${}^*\Q$. Interpret each added function by applying its
symbol to representative ground terms. Constraint (b) makes the result
independent of representatives. Every element of $R$ is represented by
a ground term, so (c) implies every universal Skolemized axiom in $R$.
Its reduct is the required model of $\mathcal T$, and (d) gives the
pole-parity condition on every element. The interpretation of fraction
fields extends the embedding to
$F_R\hookrightarrow{}^*F$.

The ambient diagram is used only in this existence argument. It is not
required to be recursive and is never queried by the finite rational
tests.
\end{proof}

We now analyze any model and embedding supplied by
Lemma~\ref{red:compactness}. All starred objects refer to its one
ambient extension. The point conditions give the additive injection
\(\eta\) of
\eqref{red:eta}--\eqref{red:additivity}.

\subsection{Odd valuations select a tested branch}
\label{red:valuation}

The height radical supplies primes of the ambient integer sort,
whereas the tests are imposed at ideals above an element \(r\in R\).
We now connect these two kinds of data. Fix an ordinary precision
\(K\ge1\). Enlarge the finite exceptional set \(S\) in the height
estimate to a finite set \(S_K\) containing:
the exceptional set for $\Phi$, the supports of all multipliers in
$\mathcal K$, the primes needed to keep the contact forms distinct
with their required unit determinants, the bad or ramified primes of
the elliptic data and $F$, and the finitely many primes excluded by
Proposition~\ref{ell:formal} at precision $K$.

Write \(M_K(s)=M_{S_K}(s)\) for the radical with these larger
exceptions. Corollary~\ref{hei:enlarge-exceptions} keeps the exponent
unchanged: if \(M(s)=M_S(s)\) is the original radical, then
\[
 M(s)\le\left(\prod_{q\in S_K\setminus S}q\right)M_K(s),
 \qquad \hgt(s)\le H_K M_K(s)^c
\]
with an ordinary constant $H_K$. Enlarge $c$ to a positive integer
once and for all if necessary.

Choose the representation axiom at precision $K$ for any $A\in R$,
and put $\widetilde s_i=(\widetilde u_i:\widetilde v)$.
All its contact forms are nonzero, so
$\widetilde s_i\notin\mathcal C$. Consider any ambient prime $q$
counted in some $M_K(\widetilde s_i)$. Here and below a radical, prime,
or valuation on starred sorts is the transfer of its ordinary definition.

Let $t=\min\{v_q(\widetilde u_i),v_q(\widetilde v)\}$.
At $q\notin S_K$, the minimum of the valuations of the five contact
forms equals $t$: all their coefficients are integral and a pair has
unit determinant. Since some contact depth is odd, these five
valuations cannot all have the same parity. In particular one of them
is odd. For its expression $\varepsilon k r$ in
\eqref{red:contact-list}, $k$ is a $q$-unit, and hence
\begin{equation}\label{red:odd-r}
                            v_q(\widetilde r)\text{ is odd}.
\end{equation}

The odd valuation links this ring element \(r\) to the ambient
prime \(q\), but it does not make the tested localization a subring of
the ambient valuation ring: \(R\) may have denominators at \(q\).
We will quotient the value group by the convex subgroup generated by
the negative valuations of nonzero ring elements. Pole parity ensures
that the odd value of \(r\) survives positively in this quotient. The
center of the resulting valuation will then select one of the two
tested branches.

\begin{lemma}\label{red:coarsening}
For the prime $q$ and element $r$ in \eqref{red:odd-r}, there is a
branch $\mathfrak p\in\{\mathfrak p_+,\mathfrak p_-\}$ of that
representation axiom and an ambient place $v_*$ of ${}^*F$ over $q$
such that, for every ordinary integer $l\ge1$ and every $D'\in F_R$,
\begin{equation}\label{red:valuation-transfer}
 D'\in r^l(R_F)_{\mathfrak p}
               \quad\Longrightarrow\quad v_*(\widetilde D')\ge l.
\end{equation}
Every point formal for this tested localization actually reduces to
$O$ for $v_*$.
\end{lemma}

\begin{proof}
Let $J$ be the convex subgroup of ${}^*\Z$ generated by the negative
values of $v_q$ on $\widetilde R\setminus\{0\}$. Thus $\gamma\in J$
if $|\gamma|$ is bounded by a finite sum of absolute values of such
negative values. This definition is external; we will use it only to
construct a valuation on the embedded ring, not as part of an algorithm
or a transferred formula.

Every valuation of a nonzero ring element that lies in $J$ is even.
If $J=0$, this is immediate.
Otherwise, given $a\ne0$ with $v_q(\widetilde a)\in J$, finite products
of ring elements having negative value give $b\in R\setminus\{0\}$
with
$v_q(\widetilde b)<-|v_q(\widetilde a)|$. Both $\widetilde b$ and
$\widetilde a\widetilde b$ have negative valuation. The soundness of
$\Phi$, transferred from Proposition~\ref{par:pole-parity}, says that
these two valuations are even. Their difference
$v_q(\widetilde a)$ is therefore even as well.

By \eqref{red:odd-r}, $v_q(\widetilde r)$ is outside $J$; it is
positive, because every negative ring value lies in $J$.
Extend $v_q$ to an ambient place $v_*$ on ${}^*F$. Outside the fixed
ramified primes, its value group is still ${}^*\Z$ with the same
normalization, by transfer of the corresponding ordinary place facts.
Coarsen it to the ordered quotient ${}^*\Z/J$, obtaining $w_*$.
This coarsened valuation is nonnegative on $\widetilde R$:
all its negative original values become zero. It is also nonnegative
on $\widetilde R_F$, since $\sqrt2$ is integral.

The pullback of the center on the embedded \(R_F\) is
\[
 \mathfrak p=\{x\in R_F:w_*(\widetilde x)>0\}.
\]
It is a proper prime ideal, and it contains $r$ because
$w_*(\widetilde r)>0$. On the other hand,
\[
 R_F/rR_F\simeq(R/rR)[X]/(X^2-2)
                         \simeq(R/rR)\times(R/rR),
\]
where the last isomorphism uses $B_r^2=2$ and the invertibility of
$2B_r$. A prime of this product is one of its two maximal ideals.
Consequently $\mathfrak p$ is exactly one of the tested branch ideals
$\mathfrak p_\pm$.

If $D'=r^l X/Y$ as in \eqref{red:localization-definition}, then
$Y\notin\mathfrak p$, so its image has $w_*$-value zero. Thus
\[
                       w_*(\widetilde D')\ge l w_*(\widetilde r).
\]
For $D'=0$, \eqref{red:valuation-transfer} is automatic. Otherwise,
if $J=0$, the displayed inequality gives
$v_*(\widetilde D')\ge l v_q(\widetilde r)\ge l$.
If $J\ne0$, its convexity implies that it contains every ordinary
integer. A positive value in the ordered quotient is represented by
an ambient integer larger than every element of $J$; the same inequality
again gives $v_*(\widetilde D')\ge l$.

Finally, a tested finite formal point has
$1/x\in r(R_F)_{\mathfrak p}$, so its actual $x$-valuation is
negative. At the good integral model this is exactly the condition
for reduction to $O$. The identity point has that reduction as well.
\end{proof}

\subsection{Index congruences at every contact prime}
\label{red:indices}

\begin{proposition}\label{red:index-congruence}
Fix an ordinary \(K\ge1\) and use the exceptional set \(S_K\) above.
Choose a representation supplied at precision \(K\) for an arbitrary
\(A\in R\), and put
\(\widetilde s_i=(\widetilde u_i:\widetilde v)\).
For every prime \(q\) of \({}^*\Z\) counted by any
\(M_K(\widetilde s_i)\), and every \(a\in R\), the element
\(\widetilde a\) is integral at \(q\) and
\begin{equation}\label{red:index-congruence-equation}
                         v_q(\widetilde a-\eta(a))\ge K.
\end{equation}
\end{proposition}

\begin{proof}
The odd-contact argument selects a factor \(r\) with odd
\(q\)-valuation. Lemma~\ref{red:coarsening} supplies one of its tested
branches and a place \(v_*\) above \(q\). At this branch the theory
supplies nonzero \(h,j_0\), fixed for all \(a\in R\).
Every tested formal point is an actual formal point at \(v_*\), for
its own curve. The same lemma sends the memberships
\eqref{red:primary-test} and \eqref{red:alternate-test} to the two
valuation comparisons of Lemma~\ref{red:local-comparison}, each at
precision \(K\) and at this same place. The primary logarithm denominator
is nonzero by its axiom and the embedding; the conjugate parameter
denominator is nonzero by \(j_0\ne0\) and the point conditions.
Finally \(q\notin S_K\) ensures \(q\notin S_E\) and \(q>K+1\).
All hypotheses of Lemma~\ref{red:local-comparison} therefore hold,
and it gives the assertion.
\end{proof}

\subsection{The height contradiction}
\label{red:height-contradiction}

\begin{proposition}\label{red:finite-test-success}
Let $f\in\Z[X_1,\ldots,X_n]$ be nonconstant. If every finite system
of rational constraints constructed in \S\ref{red:finite-tests} for
$f$ is solvable, then $f$ has a zero in $\Z^n$.
\end{proposition}

\begin{proof}
Choose the model supplied by Lemma~\ref{red:compactness}, with its
root tuple $\mathbf a$ and the additive map $\eta$ of
\eqref{red:eta}--\eqref{red:additivity}. Set
\[
             \delta=f(\eta(a_1),\ldots,\eta(a_n))\in{}^*\Z.
\]
If $\delta=0$, the ambient integer sort satisfies
$\exists\mathbf x\in\Z^n\ f(\mathbf x)=0$; by elementarity the
ordinary integers satisfy that sentence. It remains to consider
$\delta\ne0$.

Write
\begin{equation}\label{red:B-D}
 B=\max(1,|\eta(a_1)|,\ldots,|\eta(a_n)|),\qquad
 D=\max(1,\deg f).
\end{equation}
For an ordinary constant $C_f>0$, transfer of the elementary polynomial
bound gives
\begin{equation}\label{red:error-bound}
                              |\delta|\le C_f B^D.
\end{equation}
Choose an ordinary integer $K\ge cD$, where $c$ is the fixed positive
integer exponent of the height estimate. All exceptions and constants
below refer to this fixed $K$.

Figure~\ref{fig:finite-test-proof} summarizes how the chosen representations
for all \(A\in R\) use the same \(\delta\).

\suppressfloats[t]
\begin{figure}[!ht]
\centering
\begin{tikzpicture}[
  box/.style={draw=blue!45!black,fill=blue!2,rounded corners=1.5pt,
              align=center,font=\small,inner sep=5.5pt,text width=6.55cm},
  result/.style={box,text width=14.15cm,fill=blue!5},
  line/.style={-{Stealth[length=2mm]},semithick,draw=blue!45!black}]
\node[box,minimum height=1.72cm] (slopes)
  {\textbf{Vary $A\in R$}\\[3pt]
   Choose a representation at precision $K$:\\[2pt]
   $\widetilde A=\sum_{i=1}^3\widetilde u_i/\widetilde v$,
   \quad $\widetilde s_i=(\widetilde u_i:\widetilde v)$.};
\node[box,minimum height=1.72cm,right=0.65cm of slopes] (root)
  {\textbf{Keep the root tuple $\mathbf a$ fixed}\\[3pt]
   $f(\widetilde{\mathbf a})=0$, \quad
   $\delta=f(\eta(\mathbf a))\ne0$,\\[2pt]
   $|\delta|\le C_f B^D$, \quad $K\ge cD$.};

\node[box,below=0.5cm of slopes] (local)
  {For each ambient prime $q\mid M_K(\widetilde s_i)$:\\[3pt]
   one tested branch gives actual formal points
   and valuation bounds above $q$, using
   pole parity and coarsening.};

\node[box,minimum height=1.72cm,below=0.5cm of local] (comparison)
  {\textbf{Paired elliptic tests}\\[3pt]
   $\widetilde a$ is $q$-integral and
   $v_q(\widetilde a-\eta(a))\ge K$\\[2pt]
   for every $a\in R$.};
\node[box,minimum height=1.72cm,right=0.65cm of comparison] (error)
  {Apply the comparison to $a_1,\ldots,a_n$:\\[2pt]
   $v_q(\delta)\ge K$ for every such $q$.\\[3pt]
   Hence $M_K(\widetilde s_i)^K\mid\delta$
   for $i=1,2,3$.};

\node[result,minimum height=1.62cm,below=0.5cm of comparison,
  xshift=3.8cm] (height)
  {\textbf{Apply the five-point height estimate to each chosen slope}\\[3pt]
   $\displaystyle
   \hgt(\widetilde A)
     \le\sum_{i=1}^3\hgt(\widetilde s_i)+\log3
     \le 3H_K|\delta|^{c/K}+\log3
     \ll_f B\qquad(A\in R).$};

\draw[line] (slopes)--(local);
\draw[line] (local)--(comparison);
\draw[line] (comparison)--(error);
\draw[line] (root)-- node[midway,fill=white,inner sep=2pt,
  font=\small,align=center] {polynomial\\evaluation} (error);
\draw[line] (error.south)--(height.north -| error.south);
\end{tikzpicture}
\caption{The use of one nonzero $\delta$ for every $A\in R$.
All slopes shown come from the representation chosen at the fixed ordinary
precision $K$, and $M_K$ uses the fixed exceptional set for $K$.
The comparison holds for every $a\in R$, so it applies to the fixed root tuple.
The implied constant in the final bound is ordinary and independent of $A$
and $\mathbf a$.}
\label{fig:finite-test-proof}
\end{figure}


For any $A\in R$, take its representation at precision $K$ and
any prime $q$ counted in any $M_K(\widetilde s_i)$.
Proposition~\ref{red:index-congruence} applies to each $a_j$.
Both the embedded coordinates and their integer indices are
\(q\)-integral. Since \(f(\widetilde{\mathbf a})=0\) in the ambient
field and the coefficients of \(f\) are integers, their coordinatewise
congruences give
\begin{equation}\label{red:error-divisibility}
                                  v_q(\delta)\ge K.
\end{equation}
For each fixed slope, \eqref{red:error-divisibility} now holds for
every ambient prime counted by its radical. This is the internal
universal premise in the transfer of the ordinary divisibility fact:
if every prime in a squarefree product divides an integer to order at
least \(K\), then the \(K\)th power of that product divides the integer.
Using the transferred finite-product function gives
\begin{equation}\label{red:radical-divides}
                             M_K(\widetilde s_i)^K\mid\delta.
\end{equation}
The coarsenings may have been chosen externally one prime at a time;
only the universal assertion just established enters this transferred
implication. Although \(A\) and its representation were arbitrary,
the integer \(\delta\) is the same for all of them.

As $\delta\ne0$, \eqref{red:radical-divides} and
\eqref{red:error-bound} give
\[
 M_K(\widetilde s_i)^c
       \le |\delta|^{c/K}
       \le C_f^{c/K}B^{cD/K}\ll_f B.
\]
The height estimate therefore implies
\begin{equation}\label{red:slope-height}
                               \hgt(\widetilde s_i)\ll_f B.
\end{equation}
Since $v\ne0$ and $u_1+u_2+u_3=Av$, the ordinary logarithmic height
inequality for a sum of three numbers, transferred to the ambient
structure, gives
\begin{equation}\label{red:all-ring-height}
 \hgt(\widetilde A)
   =\hgt\!\left(\frac{\widetilde u_1}{\widetilde v}
             +\frac{\widetilde u_2}{\widetilde v}
             +\frac{\widetilde u_3}{\widetilde v}\right)
       \le\sum_{i=1}^3\hgt(\widetilde s_i)+\log3
       \ll_f B \qquad(A\in R).
\end{equation}
The implicit constant is ordinary and independent of $A$ and of the
root tuple. The dependence on $f$ comes only from its degree and
coefficients and the resulting fixed choice of $K$.

For each nonzero \(a_j\), formula
\eqref{red:coordinate-functions} expresses \(x(T_{a_j})\) using
the three ring elements \(x_0(a_j),x_1(a_j),x_2(a_j)\), with
\(x_2(a_j)\ne0\). Applying \eqref{red:all-ring-height} to these
three elements and using the logarithmic height inequalities over the
fixed number field \(F\) gives
\begin{equation}\label{red:elliptic-upper}
                         \hgt(x(\eta(a_j)P))\ll_f B.
\end{equation}
Canonical height on the fixed elliptic curve gives, for ordinary
integer $n\ne0$, and hence for ambient integer indices by transfer,
\begin{equation}\label{red:elliptic-lower}
             \hgt(x(nP))=2\widehat h(P)n^2+O(1),
                         \qquad\widehat h(P)>0;
\end{equation}
see \cite{Silverman2009}. The error term is uniformly bounded for
the fixed curve and point. If $B>1$, choose a maximizing nonzero
index in \eqref{red:B-D}. Equations
\eqref{red:elliptic-upper}--\eqref{red:elliptic-lower} give
\[
                         2\widehat h(P)B^2\le C B+C'
\]
for ordinary constants, so $B$ is bounded by an ordinary constant.
This is also true if $B=1$.

Every element of ${}^*\Z$ in an ordinary finite interval is an
ordinary integer, by transfer of the finite description of that
interval. Thus $\eta(a_j)=n_j\in\Z$ for each $j$.
By additivity and $\eta(1)=1$, $\eta(n_j)=n_j$; injectivity then
gives $a_j=n_j$ in $R$. The relation $f(\mathbf a)=0$ is consequently
an ordinary integer solution. In the case currently under consideration
this also contradicts $\delta\ne0$.
\end{proof}

\begin{proof}[Proof of Theorem~\ref{thm:main}]
Constant polynomials can be decided directly. For nonconstant $f$,
if every finite rational test succeeds,
Proposition~\ref{red:finite-test-success} gives an integer zero. Therefore,
when $f$ has no integer zero, some finite test must fail. Run the
finite-test search and enumeration of integer tuples in parallel,
allocating successive finite steps to each. The enumeration terminates
when a zero exists; the test search terminates when none exists.
Their answers are sound. Rational decidability would therefore give
integer decidability, contradicting the Davis--Putnam--Robinson--Matiyasevich
theorem \cite{DavisPutnamRobinson1961,Matiyasevich1970}.
\end{proof}

